\chapter{Conclusion and Future Work}

\section{The question, revisited}

This dissertation asked one question with two halves: \textit{can subtractive attention guidance be
delivered on consumer XR hardware, and when it is delivered, does it actually help human attention?}

The first half has been answered in the affirmative, with precision about what "delivered" means.
Chapter 3 showed that on the Meta Quest 3 the question reduces to a compositing-rights problem, and
mapped the three tiers of access an application actually has: a colour-only styling interface to the
OS-owned passthrough layer; overlay compositing on top of it; and full pixel control over a
lower-quality camera feed. Chapter 4 demonstrated that this narrow space is sufficient for a working
multi-mode attention-guidance system, built around an architectural inversion: the focus region
rendered as an \textit{absence} in an overlay, so that native passthrough at full quality serves the region
that matters, while ten peripheral treatments spanning the design space occupy the periphery around
it. Two of them, SignPop and Hard Dark, are carried to confirmatory evaluation. The hybrid composite at
the heart of the most capable mode, native passthrough inside the window, a shader-processed camera
feed outside it, survived a structured prior-art search (Chapter 2) with no published or shipped
precedent found. Chapter 5 specified how the system's costs are to be measured, including the
distance between live on-device detection and the offline oracle that SignPop, one of the two
evaluated modes, itself depends on for its colour pop, its local window, and its saliency boost, and
reported preliminary results from the offline detection pipeline.

The second half of the question has an interim answer, and a confirmatory one still pending. Chapter
6's pre-registered study design carries a single, properly powered primary hypothesis, that peripheral
dimming reduces the measurable cost distractors impose on a focal task, and an equivalence-bounded
safety hypothesis, that salience re-grading must not degrade peripheral awareness beyond a stated
margin. It also carries pilot gates that prevent an uninterpretable main study from launching, and a
decision table that commits the analysis to a conclusion under every outcome, once the full N = 20
sample is collected. Eleven participant pairs exist today. Chapter 6, Section 6.11 reports what they
show: a 3.65-point workstation accuracy gain, peripheral awareness that was not lost, and a detection
benefit concentrated exactly in the eccentricity band the system's own detection-gated design
predicts. Every part of that interim answer is qualified by the sample it comes from, with achieved
power of 59–64\%, well short of the pre-registered target. But "inconclusive" is still not among the
reachable outcomes. The interim data point in a specific, mechanistically coherent direction, and the
remaining collection will either confirm that direction with tighter intervals or correct it.

\section{Contributions, honestly stated}

Four contributions were claimed in Chapter 1; their standing at the time of writing is as follows.

\textbf{C1: the design space} (Chapter 3) is complete as an analytical contribution and is grounded in
implementation experience rather than speculation. Its quantitative columns (frame cost, latency,
legibility per tier) await the Chapter 5 benchmark campaign.

\textbf{C2: the reference implementation} (Chapter 4) is complete and running on target hardware: the
world-locked focus-window architecture, world-direction fusion sampling for the monocular feed,
motion-coupled suppression, and ten peripheral treatments spanning the design space, of which two,
SignPop and Hard Dark, are carried to confirmatory evaluation, together with an engineering record
that reports dead ends, screen-space sampling double vision, a deprecated surface-projection API,
shader stripping on Android, as reusable knowledge.

\textbf{C3: the technical evaluation} exists as a specified plan with preliminary results (the offline
full-resolution detection bake: 13,939 detections across 1,040 sampled frames, 99.5\% of samples with
at least one detection). The benchmark campaign itself, frame timing, latency, legibility, thermal
endurance, and the oracle-versus-live gap, remains to be run.

\textbf{C4: the pre-registered study design} is complete, ethics-mapped against the approved protocol,
and instrumented on paper down to verbatim participant scripts and exclusion rules. Data collection is
underway. Eleven of the target twenty participant pairs are in, with interim results reported in
Chapter 6, Section 6.11 and discussed in Chapter 7, Section 7.2. Recruitment continues toward the
full sample, and the confirmatory analysis and decision table of Chapter 6, Section 6.9.3 apply once
it is reached.

An examiner should read C3 and C4 as what they are: the empirical programme of this dissertation, in
a state where every analytical decision has been made and committed to in advance of the data. That
is precisely the state in which such decisions carry evidential weight.

\section{Future work}

Seven directions follow directly from the material.

\textbf{Separating the bundle.} The interim 20–30\textdegree{} finding (Chapter 6, Section 6.11) is
the clearest reason this item leads the list. A confirmed detection switches on six effects at once:
colour pop, window carve-out, saturation lift, brightness lift, contrast expansion, and darkened
surround. So nothing in this dissertation attributes the +18.6-point gain to any one of them. A
three-arm design, no filter, filter with the detection gate off, filter with it on, would need no new
code, since the gate already has a toggle. It would directly separate "dimming helped" from "passing
the signal through helped."

\textbf{Closing the oracle gap.} SignPop, one of the two evaluated modes, is detector-dependent today,
on offline, oracle-quality detection, not a live pipeline, and Chapter 5's benchmark quantifies how
far live on-device detection currently falls short of that oracle. The engineering programme this
implies, quantised detectors, region-of-interest second passes, temporally amortised inference, is
well charted. The natural next system iteration replaces SignPop's offline oracle gate with a live
detector and re-runs Block B to test whether the measured benefit survives the live-detection
accuracy and coverage gap this dissertation deliberately left unclosed.

\textbf{Transfer to optical see-through hardware.} The Tier-2 dark overlay does not transfer to OST
glasses, whose additive optics cannot darken the world. This was stated in Chapter 7 and bears
repeating as a boundary, not a defect. What does transfer is the taxonomy itself, on OST hardware
the design space collapses towards display-side dimming layers and segment attenuation, and,
critically, the human-factors findings, which concern peripheral diminishment as a perceptual
manipulation rather than any one rendering path. A replication of the Chapter 6 paradigm on a
dimming-capable OST device would be the cleanest possible test of that transfer.

\textbf{Adaptive introduction of diminishment.} An idea raised and deliberately deferred during study
design: rather than switching the effect on, a deployed system might introduce it gradually as focus
wanes. As a validation instrument a ramp confounds time with intensity and was excluded from the
study. As an \textit{adaptive-system UX question}, when and how fast should a focus aid fade in, and should
the user or the system control it, it is a well-formed research problem, seeded by the system's
existing formation animation and motion-coupled fade mechanics, and by the pre-registered
time-on-task analysis in the study, which will show whether the vignette's benefit grows as boredom
sets in.

\textbf{Dose–response.} The study tests only the extreme of the attenuation axis (Hard Dark). The
midpoint, Soft Dark, was built but never tested with participants, so the middle of the axis has no
data behind it. A discrete-levels design, off, partial,
full attenuation, counterbalanced, would establish the shape of the benefit curve and answer the
deployment question this dissertation cannot: how much diminishment buys how much focus, at what
comfort cost.

\textbf{Hour-long deployment.} The published thermal ceiling on continuous on-device camera processing
(throttling within five to ten minutes) makes the dark modes, which have no camera pipeline, the
only plausible basis for session-length deployment today. A longitudinal study of real work sessions
under one of them, with the study's workload and comfort instruments administered across days rather
than minutes, is the necessary bridge between this dissertation's controlled findings and any claim
about practice.

\textbf{Saliency-weighted diminishment.} The literature review identified the combination of channels
(desaturation with blur; dimming weighted by computed scene saliency) as unstudied in AR attention
guidance. The system's shader architecture already computes per-fragment treatment. Driving its
intensity from a saliency model of the periphery, dimming most where the periphery pulls hardest,
is a natural composition of the Tier-3 rendering path with the saliency-modulation literature, and a
candidate for the strongest version of the guidance concept.

\section{Closing}

The premise of this work was that a consumer headset can be an instrument for \textit{removing} the world as
well as adding to it, and that both halves of that proposition deserve rigour: the systems half,
where an OS-owned display layer had to be subtracted from by purely additive means, and the human
half, where "it helps you focus" had to be converted from a hope into a hypothesis that could fail.
The system exists and runs on target hardware. The hypothesis has moved past waiting: eleven
participant pairs have already tested it, and the interim direction is favourable, though the
confirmatory answer at the full N = 20 sample is still being collected. Whichever way the final data
fall, the question will have been asked properly, on the hardware people actually own.

\section*{References}

\refitem{Laghari, M. K., Shaikh, A. A., Khan, F., \& Siddiqui, A. G. (2025). Native mixed reality compositing
on Meta Quest 3: a quantitative feasibility study of ARM-based SoCs and thermal headroom. \textit{arXiv
2509.18929}. \url{https://arxiv.org/abs/2509.18929}}

\refitem{Cheng, Y., et al. (2022). Towards understanding diminished reality. \textit{Proc. CHI 2022}.}

\refitem{McLaughlin, A. C., et al. (2025). Cognitive aid design using diminished reality to support
selective attention by reducing distraction. \textit{Human Factors}, 67(9), 937–961.
doi:10.1177/00187208251325169}
